\begin{table}[H]
\centering\footnotesize
\caption{Ejemplo del plan generado para el lunes 12 y el martes 13 de octubre (caso base); cada bloque lleva su motivo}\label{tab:ejemplo}
\begin{tabularx}{\textwidth}{l l >{\raggedright\arraybackslash}p{4.6cm} >{\raggedright\arraybackslash}X}
\toprule
Día & Hora & Bloque & Motivo \\
\midrule
lunes 12/10 & 09:00--11:00 & Taller MATE ECDI (06/10) & Se entrega mañana (13/10). Por su dificultad (medio) se planea con 1 día de margen. Al ser inminente, se permiten hasta 4 h hoy en esta tarea. \\
lunes 12/10 & 11:00--12:30 & Taller CIENCBAS ASTO (06/10) & Se entrega mañana (13/10). Dificultad fácil. Al ser inminente, se permiten hasta 4 h hoy en esta tarea. \\
lunes 12/10 & 12:30--14:30 & Informe de laboratorio ISIS ARSW (09/10) & Se entrega en 4 días (16/10). Por su dificultad (difícil) se planea con 2 días de margen. Se reparte en bloques de máximo 2 h por día. \\
martes 13/10 & 07:00--09:00 & Informe de laboratorio ISIS ARSW (09/10) & Se entrega en 3 días (16/10). Por su dificultad (difícil) se planea con 2 días de margen. Se reparte en bloques de máximo 2 h por día. \\
martes 13/10 & 09:00--11:00 & Informe de laboratorio ISIS PTIA (10/10) & Se entrega en 4 días (17/10). Dificultad fácil. Se reparte en bloques de máximo 2 h por día. \\
martes 13/10 & 17:00--19:00 & Parcial 2 de ECDI & Se entrega en 7 días (20/10). Por su dificultad (difícil) se planea con 2 días de margen. Se reparte en bloques de máximo 2 h por día. \\
\bottomrule
\end{tabularx}
\end{table}
